\section{Conclusion}
We addressed the semi-supervised ITE estimation problem.
In comparison to the existing ITE estimation methods that only rely on labeled instances including treatment and outcome information, we proposed a novel semi-supervised ITE estimation method that also utilizes unlabeled instances.
The proposed method called counterfactual propagation is built on two ideas from causal inference and semi-supervised learning, namely, matching and label propagation, respectively; accordingly, we devised an efficient learning algorithm.
Experimental results using the semi-simulated real-world datasets revealed that our methods performed better in comparison to several strong baselines when the available labeled instances are limited. 
% However, at the same time, reasonable similarity design and hyper-parameter tuning remained issues.
However, this method had issues related to reasonable similarity design and hyper-parameter tuning.


%Our future work includes addressing the biased distribution of labeled instances. 
%As mentioned in the section discussing the related work
%In this paper, we did not consider such sampling biases for labeled data, but some debiasing techniques~\cite{zhou2019graph} can be successfully integrated into our framework. 

% 順番を変更
One of the possible future directions is to make use of balancing techniques such as the one used in CFR~\cite{shalit2017estimating}, which can be also naturally integrated into our model.
Our future work also includes addressing the biased distribution of labeled instances. 
As mentioned in Related work, we did not consider such sampling biases for labeled data. 
Some debiasing techniques~\cite{zhou2019graph} might also be successfully integrated into our framework. 
In addition, robustness against noisy outcomes under semi-supervised learning framework is still the open problem and will be addressed in the future.


